% === Section 9: Autonomous Vehicle Validation ===
\section{Autonomous Vehicle Validation}\label{sec:av}

This section reports the promoted autonomous-vehicle evidence row.  The current
defended AV source is the completed all-zip grouped nuPlan replay surface, not the
older Waymo-era AV lane.  The implementation still passes through legacy compatibility
schemas and file names in a few internal paths, but the public claim is
source-neutral and bounded to nuPlan replay artifacts.

\subsection{Domain Description}

The AV adapter evaluates longitudinal safety prediction: ego speed and relative
gap to the lead vehicle at 1\,s, 2\,s, and 4\,s horizons.  The ORIUS runtime sees
degraded observations caused by delay jitter, dropout, out-of-order packets,
spikes, and stale telemetry.  The defended constraints are narrow:
\begin{enumerate}[nosep]
  \item the released longitudinal action must not violate the replay-level
        TTC and predictive-entry-barrier guard;
  \item the true future target must remain inside the conformal set at the
        reported coverage level;
  \item every released action must carry a valid CertOS certificate.
\end{enumerate}
This row does not claim closed-loop road deployment, policy optimality, lane-change
planning, or complete autonomy validation.

\subsection{Dataset Boundary}

\paragraph{Source.}
The source surface is the all-zip grouped nuPlan replay build recorded under
\texttt{reports/orius\_av/nuplan\_allzip\_grouped\_runtime\_dropout\_aligned\_m15\_fulltest/}.
The processed source manifest records completed nuPlan archives and excludes partial
downloads and map-only archives before feature construction.  The runtime report
contains 1{,}531{,}104 canonical ORIUS replay rows over 37{,}344 held-out test scenarios.

\paragraph{Data splits.}
\Cref{tab:av_splits} reports the scenario counts produced by the all-zip grouped
split.  The split assignment is deterministic, and every split is non-empty.

\begin{table}[h]
\caption{AV all-zip grouped nuPlan scenario splits.}
\label{tab:av_splits}
\centering
\begin{tabular}{lr}
\toprule
\textbf{Split} & \textbf{Scenarios} \\
\midrule
Train       & 265{,}904 \\
Validation  & 37{,}115 \\
Calibration & 39{,}079 \\
Test        & 37{,}344 \\
\bottomrule
\end{tabular}
\end{table}

\subsection{Training Summary}

The all-zip grouped run trains separate conformal regressors for each target and
horizon.  These values are evidence of the replay-surface calibration stack, not a
claim of unrestricted population-level AV calibration.

\begin{table}[h]
\caption{All-zip grouped nuPlan conformal training summary.}
\label{tab:av_training}
\centering\scriptsize
\begin{tabular}{llcccc}
\toprule
\textbf{Target} & \textbf{Horizon} & \textbf{Split} &
\textbf{Base cov.} & \textbf{Wid.\ cov.} & \textbf{Wid.\ width} \\
\midrule
ego speed & 1\,s & val  & 0.899 & 0.937 & 0.47 \\
ego speed & 1\,s & test & 0.898 & 0.935 & 0.48 \\
ego speed & 2\,s & val  & 0.899 & 0.934 & 0.75 \\
ego speed & 2\,s & test & 0.897 & 0.930 & 0.75 \\
ego speed & 4\,s & val  & 0.900 & 0.931 & 1.74 \\
ego speed & 4\,s & test & 0.897 & 0.927 & 1.74 \\
\midrule
rel.\ gap & 1\,s & val  & 0.907 & 0.955 & 71.27 \\
rel.\ gap & 1\,s & test & 0.905 & 0.955 & 72.63 \\
rel.\ gap & 2\,s & val  & 0.901 & 0.951 & 80.50 \\
rel.\ gap & 2\,s & test & 0.902 & 0.953 & 81.50 \\
rel.\ gap & 4\,s & val  & 0.899 & 0.951 & 89.41 \\
rel.\ gap & 4\,s & test & 0.900 & 0.952 & 90.09 \\
\bottomrule
\end{tabular}
\end{table}

The main training signal is directional: widening improves both ego-speed and
relative-gap coverage on the all-zip grouped run.  The remaining shift-aware
relative-gap weakness is preserved in the claim boundary rather than hidden by
averaging.

\subsection{Runtime Summary}

\Cref{tab:av_runtime} compares the no-wrapper baseline, an RSS/CBF-style filter,
the trivial always-brake control, and ORIUS on the same 1{,}531{,}104 runtime steps.

\begin{table}[h]
\caption{All-zip grouped nuPlan runtime summary.}
\label{tab:av_runtime}
\centering\scriptsize
\begin{tabular}{lrrrrrr}
\toprule
\textbf{Controller} & \textbf{TSVR} & \textbf{OASG} & \textbf{CVA} &
\textbf{IR} & \textbf{Fallback} & \textbf{Audit} \\
\midrule
baseline & 0.289 & 0.023 & 1.000 & 0.000 & 0.000 & 0.000 \\
rss-cbf filter & 0.078 & 0.016 & 1.000 & 0.387 & 0.000 & 0.000 \\
always brake & 0.000 & 0.000 & 1.000 & 1.000 & 1.000 & 0.000 \\
orius & 0.000163 & 0.000157 & 1.000 & 0.500 & 0.174 & 1.000 \\
\bottomrule
\end{tabular}
\end{table}

ORIUS reduces replay-level true-state violations by 99.94\% relative to the baseline
while preserving an auditable release chain.  The always-brake control reaches zero
TSVR, but does so by activating fallback on every step and carrying no certificate
audit surface.  ORIUS therefore occupies the defended point in this run: epsilon-closed
replay release with non-trivial useful work and complete governance metadata.

\subsection{Fault Family Coverage}

\Cref{tab:av_fault} reports ORIUS coverage by fault family and target.  The table is
deliberately not summarized as uniformly perfect coverage: relative gap remains the
harder target and is kept inside the bounded AV claim.

\begin{table}[h]
\caption{ORIUS fault-family coverage on the all-zip grouped nuPlan replay row.}
\label{tab:av_fault}
\centering\scriptsize
\begin{tabular}{llrr}
\toprule
\textbf{Fault family} & \textbf{Target} & \textbf{Coverage} & \textbf{Mean width} \\
\midrule
delay\_jitter  & ego speed & 0.988 & 0.76 \\
delay\_jitter  & rel.\ gap  & 0.930 & 117.63 \\
drift\_combo   & ego speed & 0.983 & 0.70 \\
drift\_combo   & rel.\ gap  & 0.919 & 102.17 \\
dropout        & ego speed & 0.998 & 1.21 \\
dropout        & rel.\ gap  & 0.936 & 193.29 \\
out\_of\_order & ego speed & 0.983 & 0.64 \\
out\_of\_order & rel.\ gap  & 0.923 & 102.13 \\
spikes         & ego speed & 0.998 & 1.18 \\
spikes         & rel.\ gap  & 0.939 & 191.47 \\
stale          & ego speed & 0.982 & 0.62 \\
stale          & rel.\ gap  & 0.914 & 93.81 \\
\bottomrule
\end{tabular}
\end{table}

\subsection{Shift-Aware Runtime}

\Cref{tab:shift_aware_av} reports the shift-aware runtime table emitted by the
all-zip grouped nuPlan run.

\begin{table}[h]
\caption{Shift-aware statistics for the all-zip grouped nuPlan AV row.}
\label{tab:shift_aware_av}
\centering\small
\begin{tabular}{lrrrrl}
\toprule
\textbf{Target} & \textbf{Steps} & \textbf{Groups} & \textbf{Max gap} &
\textbf{Validity} & \textbf{Status} \\
\midrule
\texttt{ego\_speed\_mps} & 1{,}531{,}104 & 2355 & 0.900 & 0.844 & nominal \\
\texttt{relative\_gap\_m} & 1{,}531{,}104 & 2480 & 0.900 & 0.704 & watch \\
\bottomrule
\end{tabular}
\end{table}

The ego-speed target remains nominal.  Relative gap enters watch status, which is
consistent with the weaker spike and short-horizon relative-gap cells in
\cref{tab:av_fault,tab:av_training}.  This is a bounded replay result, not a
deployment certificate.

\subsection{Certificate Chain}

CertOS re-verification checks 1{,}531{,}104 issued ORIUS certificates.  The
certificate-valid release rate is 0.999924, all required payload fields are present
on the promoted certificate surface, and audit completeness is 1.000.  The
certificate result establishes release accountability on the replay surface; it does
not convert the bounded nuPlan run into a full autonomous-driving field claim.

\subsection{Feature Engineering Detail}

Each nuPlan replay scenario is projected into the source-neutral ORIUS AV feature
surface used by the runtime wrapper.  The compatibility schema retains historical
field names such as \texttt{shard\_id}, but scenario identifiers include archive
identity so multiple nuPlan archives cannot collide.

\begin{table}[h]
\caption{AV feature surface used by the all-zip grouped nuPlan row.}
\label{tab:av_features}
\centering\small
\begin{tabular}{llr}
\toprule
\textbf{Category} & \textbf{Features} & \textbf{Count} \\
\midrule
Ego kinematics & speed, acceleration, heading, yaw rate & 4 \\
Lead context & relative gap, relative speed, TTC & 3 \\
Scene context & nearby-agent count, density, speed statistics & 4 \\
Map/context metadata & route and lane-context indicators & 3 \\
Fault indicators & tested degradation-family indicators & 5 \\
Temporal metadata & horizon, step, split, source identity & 6 \\
Conformal/runtime & interval, reliability, certificate fields & 13 \\
\midrule
\textbf{Total} & source-neutral runtime feature surface & \textbf{38} \\
\bottomrule
\end{tabular}
\end{table}

\subsection{Visual Evidence}

\begin{figure}[h]
\centering
\includegraphics[width=0.90\textwidth]{training_interval_widths}
\caption{All-zip grouped nuPlan conformal interval widths by target and horizon.}
\label{fig:av_training_widths}
\end{figure}

\begin{figure}[h]
\centering
\includegraphics[width=0.90\textwidth]{shift_aware_runtime}
\caption{All-zip grouped nuPlan shift-aware runtime dashboard.  Relative gap remains in
watch status and is treated as a limitation of the AV row.}
\label{fig:av_shift_aware}
\end{figure}

\begin{figure}[h]
\centering
\includegraphics[width=0.85\textwidth]{fault_family_coverage}
\caption{All-zip grouped nuPlan fault-family coverage heatmap for the ORIUS controller.}
\label{fig:av_fault_coverage}
\end{figure}

\begin{figure}[h]
\centering
\includegraphics[width=0.90\textwidth]{runtime_metrics}
\caption{All-zip grouped nuPlan runtime metrics for baseline comparators and ORIUS.}
\label{fig:av_runtime_metrics}
\end{figure}

\subsection{Summary}

The AV row now says exactly what the artifacts support: on the all-zip grouped
completed nuPlan replay surface, ORIUS reduces baseline TSVR from 0.2893 to
0.000163, emits 1{,}531{,}104 certificates with a 0.999924 certificate-valid release
rate, and exposes the relative-gap shift weakness as an open limitation.  The row
validates the ORIUS replay interface and governance chain under longitudinal AV
semantics; it does not claim full closed-loop autonomy.
